% SPDX-License-Identifier: GPL-3.0-or-later
% Copyright (C) 2026 Alessandro Greco
\section*{Глава \ref{cap:28-opsec}: ошибки операционной безопасности}

\begin{soluzione}{es:28-livello}
Трафик вне Tor, документы, открытые при подключении к сети, и утечки DNS относятся к уровню приложений; меры защиты — Tor Browser, Tails или Whonix, \opt{SafeSocks} и открытие файлов без сетевого подключения. Метаданные, вход в учётные записи и смешение идентичностей относятся к уровню содержимого; нужны инструменты очистки метаданных и раздельные сеансы. Стиль письма, время публикации и использование Tor в небольшом множестве анонимности относятся к поведению; здесь помогают изменение привычек, мосты и задержка перед публикацией.
\end{soluzione}

\begin{soluzione}{es:28-dns}
Запрос DNS уходит с компьютера пользователя на настроенный в системе сервер разрешения имён, обычно принадлежащий провайдеру связи, и содержит доменное имя сайта. Провайдер и любой наблюдатель на этом участке узнают, какой сайт пользователь намеревается посетить, даже если само соединение затем пойдёт через Tor.
\end{soluzione}

\begin{soluzione}{es:28-evento}
Через Stem подпишитесь на \texttt{EventType.STATUS\_CLIENT}. Запрос с одним лишь IP-адресом вызывает событие уровня \texttt{WARN} с действием \texttt{DANGEROUS\_SOCKS} и аргументами \texttt{PROTOCOL} и \texttt{ADDRESS}. Графический интерфейс мог бы предупреждать с указанием приложения и адреса назначения, предлагая исправить его настройки SOCKS.
\end{soluzione}

\begin{soluzione}{es:28-insieme}
При использовании obfs4 наблюдатель увидел бы зашифрованное соединение без узнаваемой структуры с обычным на вид адресом и не смог бы уверенно определить, что оно связано с Tor. Остались бы другие признаки: объём и время передачи, необычная длительность соединения, похожие подключения, а также содержание сообщений и мотив, сужающие круг подозреваемых до тех, кому предстояло сдавать этот экзамен.
\end{soluzione}

\begin{soluzione}{es:28-modello}
Упражнение с открытым ответом. Хороший анализ следует структуре \cref{cap:04-modello-minaccia}: противники и их возможности, что каждый способен наблюдать на уровнях \cref{fig:livelli}, выбранные средства защиты и обоснование их выбора, а также явный перечень неустранённых рисков — например, корреляция трафика наблюдателем обоих концов или изъятие устройств.
\end{soluzione}

\section*{Глава \ref{cap:29-ricerca}: академические исследования Tor}

\begin{soluzione}{es:29-entropia}
$H = \tfrac12 \cdot 1 + \tfrac14 \cdot 2 + 2 \cdot \tfrac18 \cdot 3 = 1{,}75$ бита; $H_M = \log_2 4 = 2$ бита; $d = 0{,}875$. Эквивалентное число равновероятных пользователей: $2^{1{,}75} \approx 3{,}4$.
\end{soluzione}

\begin{soluzione}{es:29-metodo}
(a) Моделирование сети с помощью Shadow или, в небольшом масштабе, эмуляция через Chutney: необходимо поведение настоящего tor, но привлекать реальных пользователей не нужно. (b) Измерение публичной сети по опубликованным агрегированным данным, как в Tor Metrics (\cref{sec:metrics-utenti}). (c) Моделирование поведения по образцу TorPS с использованием настоящих консенсусов за длительный период, при необходимости с последующим теоретическим анализом.
\end{soluzione}

\begin{soluzione}{es:29-durata}
При \texttt{guard\_lifetime} в 30 дней (фактический срок 18--30 дней из-за случайного сдвига даты в прошлое) за 180 дней происходит примерно семь или восемь выборов входного узла, а моделирование даёт 32,1\% скомпрометированных клиентов. При стандартных 60 днях (фактически 48--60) получается примерно три выбора и 18,4\%; при 180 днях — один или два выбора и 8,6\%. Приближённо после $k$ независимых выборов доля скомпрометированных приближается к $1-(1-q)^k$ при $q = 5\%$: для $k\approx 7{,}5$ это около 31\%, что соответствует результату для 30 дней. Чем дольше служит входной узел, тем реже выборы и тем меньше новых клиентов подвергаются риску; однако клиент, выбравший враждебный узел, остаётся под наблюдением дольше.
\end{soluzione}

\begin{soluzione}{es:29-asimmetria}
Ни один из двух противников не компрометирует клиентов согласно критерию лаборатории, для которого необходимо видеть оба конца соединения. Однако получаемые ими сведения различаются: выходной ретранслятор видит назначения и трафик, не защищённый прикладным шифрованием, но не знает личности клиентов; входной ретранслятор видит IP-адреса выбравших его клиентов, но не знает их назначения и может пытаться определить посещаемые сайты по отпечаткам трафика (\cref{sec:att-wf}).
\end{soluzione}

\begin{soluzione}{es:29-torps}
Исследовательское упражнение. Повторное использование цепей сокращает число новых выборов за день и замедляет кривую компрометации при отсутствии постоянных входных узлов. Ограничение /16 обычно слабо влияет на результат, кроме случаев, когда враждебные ретрансляторы сосредоточены в одной сети. Противник на уровне автономной системы меняет модель: он наблюдает трафик без управления ретрансляторами, поэтому нужны данные о топологии Интернета, как в \cite{johnson2013users}.
\end{soluzione}

\section*{Глава \ref{cap:30-arti}: Arti}

\begin{soluzione}{es:30-crate}
Глава~5: \texttt{tor-llcrypto}; глава~7: \texttt{tor-proto} (каналы) и \texttt{tor-chanmgr}; глава~8: \texttt{tor-cell}; главы~9--11: \texttt{tor-proto} (цепи, шифрование, потоки и управление потоком); глава~12: \texttt{tor-netdoc}, \texttt{tor-dirmgr}, \texttt{tor-dirclient}; глава~13: \texttt{tor-circmgr}, \texttt{tor-guardmgr}, \texttt{tor-relay-selection}; глава~14: \texttt{tor-hsclient}, \texttt{tor-hsservice}, \texttt{tor-hscrypto}; глава~15: \texttt{equix} и \texttt{hashx} для доказательства выполненной работы; глава~16: \texttt{tor-ptmgr}.
\end{soluzione}

\begin{soluzione}{es:30-obsoleto}
Преимущества: меньше кода и поверхность атаки, отсутствие особых случаев (например, двухбайтовых CircID) и необходимости поддерживать слабую устаревшую криптографию. Недостаток: Arti несовместим с очень старыми ретрансляторами. Однако публичная сеть Tor сама уже исключает нерекомендуемые версии, поэтому это важно главным образом для частных сетей и давно не обновлявшихся установок.
\end{soluzione}

\begin{soluzione}{es:30-cella}
\texttt{decode\_cell} читает команду из пятого байта и определяет, имеет ли ячейка переменную длину. В этом случае два следующих байта задают длину полезных данных, а общая длина составляет $7 +$ указанная длина; иначе ячейка занимает 514 байт. Если байтов пока недостаточно, функция возвращает \texttt{Ok(None)}, а не ошибку: ячейка не повреждена, а лишь ещё не получена целиком. Вызывающая сторона повторяет попытку при поступлении следующих байтов. Так же должен работать потоковый декодер из Лаборатории~2.
\end{soluzione}

\begin{soluzione}{es:30-isolamento}
Да: SOCKS-прокси Arti по умолчанию изолирует потоки с разными учётными данными SOCKS и разделяет разные прослушивающие сокеты. Проверить это можно, наблюдая цепи со стороны ретрансляторов (например, через ControlPort выходного узла) или в подробных журналах Arti. При использовании Arti как библиотеки явную изоляцию обеспечивают \texttt{TorClient::isolated\_client} и токены изоляции \texttt{StreamPrefs}. На момент Arti 2.6.0 эквивалентную настройку для отдельных портов файл \file{cfg.rs} помечает как будущую возможность.
\end{soluzione}

\begin{soluzione}{es:30-libreria}
Упражнение по программированию: создайте \texttt{TorClientConfig} из TOML-файла лаборатории, запустите клиент методом \texttt{TorClient::create\_bootstrapped} и откройте поток с помощью \texttt{connect(("127.0.0.1", 8080))}, затем передайте по нему HTTP-запрос. Конфигурация должна разрешать локальные адреса, как файл, создаваемый Лабораторией~20.
\end{soluzione}

\section*{Глава \ref{cap:31-futuro}: открытые проблемы}

\begin{soluzione}{es:31-stato}
Примеры на дату, принятую в главе: \proposta{332} (ntor-v3) имеет состояние \emph{Closed}: реализована и включена в спецификации. \proposta{324} (управление перегрузкой) — \emph{Finished}: реализована, но текст ещё не целиком перенесён в спецификации. \proposta{359} (CGO) — \emph{Open}: реализация уже присутствует в коде, однако предложение ещё обсуждается. Состояние можно узнать в заголовке каждой инициативы и в указателе по состояниям репозитория.
\end{soluzione}

\begin{soluzione}{es:31-celle}
В традиционном формате расчёт по размерам даёт $\lceil 1184/498 \rceil = 3$ ячейки для открытого ключа ML-KEM-768 длиной 1184 байта и ещё $\lceil 1088/498 \rceil = 3$ для шифртекста длиной 1088 байт: всего шесть. Однако действующий протокол даже не позволяет так фрагментировать одно ретрансляционное сообщение. \proposta{340} предлагает необходимое обрамление: одно сообщение может занимать несколько ячеек, а несколько сообщений — помещаться в одной. Тело ретрансляционной ячейки в новом формате содержит 493 байта, а не 498: 16 байт отведены под тег аутентификации, $509 - 16 = 493$, как объясняется в \cref{cap:08-celle}. Первый фрагмент содержит не более 490 байт полезных данных, поскольку в нём есть и заголовок обрамления. Пока предложение ограничивает обычные тела сообщений 493 байтами; чтобы передавать более длинный постквантовый \texttt{EXTEND2}, потребуется новая версия подпротокола \texttt{Relay}, снимающая ограничение. Даже при таком расширении открытый ключ займёт три ячейки, а шифртекст — ещё три. Значения передаются в \emph{противоположных} направлениях (клиент — ретранслятор и обратно), поэтому объединить их нельзя. Всего остаётся шесть ячеек.
\end{soluzione}

\begin{soluzione}{es:31-epistemico}
Если бы каждый клиент загружал собственное подмножество ретрансляторов, противник, знающий или влияющий на подмножество конкретного пользователя, мог бы распознавать его цепи: цепь через узлы, известные лишь этому пользователю, раскрывала бы, кто её создал. Даже без предварительных сведений наблюдение за тем, какие ретрансляторы используются со временем, сужало бы множество возможных клиентов.
\end{soluzione}

\begin{soluzione}{es:31-conflux}
Отберите цепи с назначением \texttt{CONFLUX\_LINKED} и сгруппируйте по \texttt{CONFLUX\_ID} (\cref{cap:18-stem}), после чего сравните маршруты каждой пары. Выходной ретранслятор всегда один и тот же. Другие позиции в лабораторной сети часто тоже совпадают: например, две ветви \texttt{[test007e, test003r, test006e]} и \texttt{[test003r, test007e, test006e]} состоят из тех же трёх ретрансляторов в другом порядке. Это не ошибка tor: в лаборатории мало подходящих узлов, а \proposta{329} выражает пожелание через \texttt{SHOULD NOT}, а не строгий запрет. В публичной сети с тысячами ретрансляторов такое совпадение гораздо реже.
\end{soluzione}

\begin{soluzione}{es:31-proposta}
Упражнение с открытым ответом. Образец структуры задаёт \proposta{001}: заголовок с номером, названием, автором, датой и состоянием, затем мотивация, описание решения, анализ безопасности и совместимости, план реализации и внедрения.
\end{soluzione}

\section*{Приложение \ref{cap:B-crittografia}: дополнительная криптография}

\begin{soluzione}{es:B-bit}
Из 256 бит пять фиксируются (три младших, старший и предпоследний), поэтому остаётся $2^{251}$ значений вместо $2^{256}$. Это не ослабляет ключ: безопасность кривой всё равно составляет примерно 128 бит.
\end{soluzione}

\begin{soluzione}{es:B-inversa}
Из $y = (u-1)/(u+1)$ следует $u = (1+y)/(1-y) \bmod p$. В Python значение $y$ читают из 32 байтов в порядке от младшего к старшему (little-endian), очистив бит знака; затем вычисляют \texttt{(1 + y) * pow(1 - y, P - 2, P) \% P} и записывают результат в 32 байта в том же порядке. Функция \texttt{montgomery\_to\_edwards} с исходным битом знака $x$ восстанавливает исходный ключ Эдвардса (при $y \neq 1$).
\end{soluzione}

\begin{soluzione}{es:B-digest}
Случайно изменённая ячейка проходит проверку примерно с вероятностью $2^{-16} \cdot 2^{-32} = 2^{-48}$: два нулевых байта \texttt{recognized} и четыре верных байта дайджеста. Однако при атаке маркировки не требуется, чтобы изменённую ячейку \emph{приняли}. Сообщник, видящий её, или конечный узел, отвергающий её, узнаёт об изменении по последствиям — например, закрытию цепи или предсказуемому содержимому — и этого достаточно для связывания двух участков цепи.
\end{soluzione}
